\chapter{Cross Spectra, Lags, and Coherence}

Cross-spectral analysis in \ghats{} starts from two compatible
\texttt{.fft} files.  The two files should normally be produced from the
same event data, with the same time bin, transform length, GTI selection,
and start time, changing only the channel or energy selection.  The routines
then average the complex cross spectrum and return phase or time lags,
coherence, and, in the newer interface, the real and imaginary parts of the
cross spectrum.

The routines described here are:

\begin{itemize}
\item \routine{gh_cross}, the older lag/coherence reader;
\item \routine{gh_cross_re_im_new}, the lower-level routine that also
  returns the real and imaginary parts;
\item \routine{gh_cs}, the recommended wrapper for most current work;
\item \routine{gh_plot_lag} and \routine{gh_plot_coh}, described in the
  previous chapter, for quick-look plots.
\end{itemize}

\section{Lag Sign and Basic Convention}

The first filename is the reference band and the second filename is the
subject band:

\begin{ghatsconsole}
Ghats> file_ref = 'soft.fft'
Ghats> file_sub = 'hard.fft'
\end{ghatsconsole}

Positive lags mean that the second time series lags the first.  By default
the lag array is a phase lag in radians.  If \keyword{TLAG} is set, the
routine divides by \(2\pi f\) and returns time lags in seconds:

\begin{ghatsconsole}
Ghats> gh_cs, file_ref, file_sub, freq, lag, lag_err, coh, coh_err, $
       -100, real, real_err, imag, imag_err, modcs, modcs_err, mw, qf

Ghats> gh_cs, file_ref, file_sub, freq, tlag, tlag_err, coh, coh_err, $
       -100, real, real_err, imag, imag_err, modcs, modcs_err, mw, qf, /tlag
\end{ghatsconsole}

The rebin factor \param{irf} follows the usual GHATS convention.  A positive
integer rebins linearly by that number of bins.  A negative value requests
logarithmic rebinning; \texttt{-100} is a common exploratory choice.

\section{The Recommended Interface: \routine{gh_cs}}

\routine{gh_cs} is the current high-level cross-spectral command.  It wraps
\routine{gh_cross_re_im_new} and, when \param{POISSON=[f1,f2]} is given,
uses a preliminary pass to estimate the mean real cross-spectrum contribution
in the requested high-frequency range.  It then subtracts that constant real
component in the final pass.

\begin{ghatsconsole}
Ghats> gh_cs, file_ref, file_sub, freq, lag, lag_err, coh, coh_err, $
       -100, real, real_err, imag, imag_err, modcs, modcs_err, mw, qf, $
       poisson=[400,800], rms1=0.001, rms2=0.001, /lowcoh
\end{ghatsconsole}

The output arrays are:

\begin{description}
\item[\param{freq}] Rebinned Fourier frequencies, in Hz.
\item[\param{lag}] Phase lags in radians, or time lags in seconds with
  \keyword{TLAG}.
\item[\param{lag_err}] Lag errors.
\item[\param{coh}] Coherence spectrum.
\item[\param{coh_err}] Coherence errors.
\item[\param{real}, \param{imag}] Real and imaginary parts of the averaged
  cross spectrum.
\item[\param{real_err}, \param{imag_err}] Errors on the real and imaginary
  parts.
\item[\param{modcs}] Cross-spectrum modulus,
  \(\sqrt{\mathrm{Re}(C)^2+\mathrm{Im}(C)^2}\).
\item[\param{modcs_err}] Error on the modulus, propagated using the empirical
  covariance between the real and imaginary parts.
\item[\param{mw}] Number of averaged cross vectors contributing to each
  rebinned bin.  It is the number of selected transforms times the number of
  native frequencies in the rebinned bin.
\item[\param{qf}] Quality flag for the \keyword{LOWCOH} calculation.  A value
  of 0 marks the standard high-power/high-coherence case, 2 marks the
  low-coherence case, and 1 marks bins where the implemented conditions do not
  allow a useful intrinsic-coherence estimate.
\end{description}

\begin{ghatsnote}
Use \routine{gh_cs} unless you specifically need to control the two passes
yourself.  It is less error-prone than manually measuring and subtracting the
constant real component.
\end{ghatsnote}

\section{Compatible FFT Files and Selections}

The routines check that the two FFT files are compatible.  They require the
same time resolution, Fourier length, start time, and number of transforms.
The observatory and target keywords must also agree.  The current
\routine{gh_cross_re_im_new} code warns, but does not stop, if the
instrument strings differ.

The selection keywords are the same style as in the PDS/FFT readers:

\begin{description}
\item[\param{INDEX=[i1,i2]}] Select an internal transform-index range.
\item[\param{TIME=[t1,t2]}] Select transforms by time range.
\item[\param{RATE=[r1,r2]}] Select transforms by stored count-rate range.
\item[\param{SEL=sel}] Select an explicit array of transform indices.
\end{description}

The selected transform list should match between the two FFT files.  If the
files were produced from the same data with the same production parameters,
the transform numbers should normally correspond one-to-one.

\section{Noise Prescriptions for Coherence}

\routine{gh_cs} and \routine{gh_cross_re_im_new} distinguish raw measured
coherence from intrinsic coherence.  The choice is explicit.

\subsection{\param{POISSON=[f1,f2]}}

\param{POISSON=[f1,f2]} estimates the PDS noise levels from a frequency range
and uses them in the intrinsic-coherence denominator.  In \routine{gh_cs},
the same frequency range is also used to estimate the constant real
cross-spectrum contribution that is subtracted from the real part:

\begin{ghatsconsole}
Ghats> gh_cs, file_ref, file_sub, freq, lag, lag_err, coh, coh_err, $
       -100, real, real_err, imag, imag_err, modcs, modcs_err, mw, qf, $
       poisson=[400,800]
\end{ghatsconsole}

The requested range is clipped internally to available non-Nyquist
frequencies when necessary.  The Nyquist bin is excluded from the noise
estimate and from the real-part correction.

\subsection{\param{MANUAL_POI=[p1,p2]}}

\param{MANUAL_POI=[p1,p2]} supplies fixed PDS noise levels for the two input
bands:

\begin{ghatsconsole}
Ghats> gh_cs, file_ref, file_sub, freq, lag, lag_err, coh, coh_err, $
       -100, real, real_err, imag, imag_err, modcs, modcs_err, mw, qf, $
       manual_poi=[2.01,2.03]
\end{ghatsconsole}

The two constants are interpreted in the input-file PDS units, before any
\param{RMS1}/\param{RMS2} conversion.  This option affects the coherence
calculation only; it does not apply a real-part correction to the cross
spectrum.

\subsection{\keyword{RAWCOH}}

\keyword{RAWCOH} returns measured coherence.  The observed PDS are used in
the denominator and no Poisson or noise correction is applied:

\begin{ghatsconsole}
Ghats> gh_cs, file_ref, file_sub, freq, lag, lag_err, coh, coh_err, $
       -100, real, real_err, imag, imag_err, modcs, modcs_err, mw, qf, /rawcoh
\end{ghatsconsole}

This is mutually exclusive with \param{POISSON}, \param{MANUAL_POI}, and
\keyword{LOWCOH}.

\begin{ghatsnote}
If none of \param{POISSON}, \param{MANUAL_POI}, or \keyword{RAWCOH} is
given, RXTE/PCA data use the historical Zhang-style Poisson correction.
For other instruments the routine returns raw coherence and prints a warning
suggesting \param{POISSON=[f1,f2]} or \param{MANUAL_POI=[p1,p2]} for
intrinsic coherence.
\end{ghatsnote}

\section{The \keyword{LOWCOH} Option}

\keyword{LOWCOH} changes the coherence-error calculation.  It uses the
Vaughan--Nowak low-measured-coherence formalism where the implemented
conditions indicate that the ordinary high-coherence approximation is not the
appropriate one.  The returned \param{qf} array records which path was used
for each frequency bin.

Some bins can fall outside the implemented useful range.  In that case the
routine marks the corresponding internal error term as undefined.  In IDL this
can print:

\begin{ghatsconsole}
% Program caused arithmetic error: Floating illegal operand
\end{ghatsconsole}

When it appears immediately after a \keyword{LOWCOH} run, this message is
normally the result of marking those undefined bins.  Other floating-point
messages should still be investigated.

\section{RMS Units for the Cross Spectrum}

Without \param{RMS1} and \param{RMS2}, the returned real and imaginary parts
of the cross spectrum are in the native Leahy-like units used internally by
the FFT products.  If both background rates are supplied, the routine converts
\param{real}, \param{imag}, their errors, and the modulus to rms units:

\begin{ghatsconsole}
Ghats> gh_cs, file_ref, file_sub, freq, lag, lag_err, coh, coh_err, $
       -100, real, real_err, imag, imag_err, modcs, modcs_err, mw, qf, $
       poisson=[400,800], rms1=0.001, rms2=0.001
\end{ghatsconsole}

Both background rates must be supplied together.  The first value corresponds
to the first FFT file and the second value to the second FFT file.

\section{Rotating the Cross Vector}

The \keyword{ROTATE} keyword rotates the returned real and imaginary parts of
the cross spectrum by \(\pi/4\):

\begin{ghatsconsole}
Ghats> gh_cs, file_ref, file_sub, freq, lag, lag_err, coh, coh_err, $
       -100, real, real_err, imag, imag_err, modcs, modcs_err, mw, qf, $
       poisson=[400,800], /rotate
\end{ghatsconsole}

This can be useful for fitting the real and imaginary parts when the
unrotated imaginary part is small or changes sign.  The rotation does not
change the modulus, and it does not change the returned lag or coherence
arrays.  If the rotated real and imaginary parts are fitted later, the same
rotation must be included when comparing the model to those fitted
components.

\section{Using \routine{gh_cross_re_im_new} Directly}

\routine{gh_cross_re_im_new} exposes the lower-level calculation:

\begin{ghatsconsole}
Ghats> gh_cross_re_im_new, file_ref, file_sub, freq, lag, lag_err, $
       coh, coh_err, -100, real, real_err, imag, imag_err, $
       modcs, modcs_err, mw, qf
\end{ghatsconsole}

It is useful when the real-part correction has to be controlled manually.
For example:

\begin{ghatsconsole}
Ghats> gh_cross_re_im_new, file_ref, file_sub, freq, lag, lag_err, $
       coh, coh_err, -100, real, real_err, imag, imag_err, $
       modcs, modcs_err, mw, qf
Ghats> w = where(freq ge 400 and freq le 800)
Ghats> poivalue = mean(real[w])
Ghats> gh_cross_re_im_new, file_ref, file_sub, freq, lag, lag_err, $
       coh, coh_err, -100, real, real_err, imag, imag_err, $
       modcs, modcs_err, mw, qf, poivalue=poivalue
\end{ghatsconsole}

For ordinary use, the \param{POISSON=[f1,f2]} option in \routine{gh_cs}
performs this two-pass procedure automatically.

\section{The Older \routine{gh_cross} Routine}

\routine{gh_cross} reads two compatible FFT files and returns only frequency,
lag, lag error, coherence, and coherence error:

\begin{ghatsconsole}
Ghats> gh_cross, file_ref, file_sub, freq, lag, lag_err, coh, coh_err, -100
\end{ghatsconsole}

It is retained for older workflows.  New analyses that need the real and
imaginary parts, rms conversion, rotation, the modulus, or the current raw and
intrinsic coherence choices should use \routine{gh_cs} or
\routine{gh_cross_re_im_new}.

\section{Writing Cross-Spectral Products for XSPEC}

The arrays returned by \routine{gh_cs} can be written as simple OGIP-style
PHA/RMF pairs using \routine{gh_xspec}:

\begin{ghatsconsole}
Ghats> gh_xspec, freq, real, real_err, 'real_soft_hard'
Ghats> gh_xspec, freq, imag, imag_err, 'imag_soft_hard'
Ghats> gh_xspec, freq, lag,  lag_err,  'plag_soft_hard'
Ghats> gh_xspec, freq, coh,  coh_err,  'cohe_soft_hard'
Ghats> gh_xspec, freq, modcs, modcs_err, 'mod_soft_hard'
\end{ghatsconsole}

This is a file-format conversion step.  It does not refit, rebin, or
renormalize the arrays; it writes the values and errors that were passed to
it.

\section{A Compact Workflow}

A typical current workflow is:

\begin{ghatsconsole}
Ghats> gh_nicer, 'event.evt', 'FFT', [30,200], 10000, 32768, 'soft.fft'
Ghats> gh_nicer, 'event.evt', 'FFT', [201,1200], 10000, 32768, 'hard.fft'

Ghats> gh_cs, 'soft.fft', 'hard.fft', freq, lag, lag_err, coh, coh_err, $
       -100, real, real_err, imag, imag_err, modcs, modcs_err, mw, qf, $
       poisson=[400,800], rms1=0.001, rms2=0.001, /lowcoh

Ghats> gh_plot_lag, freq, lag, lag_err, /lin
Ghats> gh_plot_coh, freq, coh, coh_err
\end{ghatsconsole}

If the reference band includes the subject band, a real-part correction is
especially important because overlapping photons introduce partial
correlation.  In that case \routine{gh_cs} with \param{POISSON=[f1,f2]} is
the preferred starting point.
